% Part: lambda-calculus
% Chapter: syntax
% Section: alpha

\documentclass[../../../include/open-logic-section]{subfiles}

\begin{document}

\olfileid{lam}{syn}{alp}
\olsection{$\alpha$-روپ بدلنا}

$\lambd[x][x]$ تے $\lambd[y][y]$ دا آپس وچ کیہ تعلق اے؟
دوویں عین اوہی قدر واپس دین والا فنکشن ظاہر کردے نیں۔ ظاہر اے کہ
نحو دے لحاظ نال ایہ وکھرے ٹرم نیں۔ ایہناں وچ فرق بس پابند
متغیر دے ناں دا اے۔ اک ٹرم دوجے وچ پابند متغیر دا ناں بدل کے
ملدا اے۔ ایس نوں \emph{$\alpha$-روپ بدلنا} آکھدے آں۔

\begin{defn}[پابند متغیر دا ناں بدلنا، $\aconvone$]
  جے ٹرم $M$ وچ $\lambd[x][N]$ دی اک آمد اے، $y \notin
  \FV{N}$ اے، تے $\Subst{N}{y}{x}$ معین اے، تاں ایس آمد دی تھاں
  \begin{equation*}
    \lambd[y][\Subst{N}{y}{x}]
  \end{equation*}
  رکھ کے $M'$ بناؤن نوں \emph{پابند متغیر دا ناں بدلنا} آکھدے
  آں۔ اینوں $M \redone[\alpha] M'$ لکھدے آں۔
\end{defn}

\begin{defn}[تعلق دی موافقت]
  ٹرماں اُتے تعلق $R$ نوں \emph{موافق} آکھدے آں جے ایہ
  ہِٹھلیاں شرطاں پوریاں کردا اے۔
  \begin{enumerate}
  \item جے $R N N'$، تاں $R \lambd[x][N] \lambd[x][N']$
  \item جے $R P P'$، تاں $R (PQ) (P'Q)$
  \item جے $R Q Q'$، تاں $R (PQ) (PQ')$
  \end{enumerate}
\end{defn}

سو تعریف نوں ہور طرح بیان کریے۔
\begin{defn}[پابند متغیر دا ناں بدلنا، $\aconvone$]
  \emph{پابند متغیر دا ناں بدلنا} ($\redone[\alpha]$) ٹرماں
  اُتے سبھ توں چھوٹا موافق تعلق اے جیہڑا ہِٹھلی شرط پوری کردا اے۔
  \begin{align*}
    &\lambd[x][N] \redone[\alpha] \lambd[y][\Subst{N}{y}{x}] && \text{جے
      $x \neq y$، $y \notin \FV{N}$} \\
    & &&\text{تے $\Subst{N}{y}{x}$ معین اے}
  \end{align*}
\end{defn}

ایتھے «سبھ توں چھوٹا» دا مطلب اے کہ تعلق وچ صرف اوہ جوڑے
نیں جیہڑے موافقت تے ودھیک شرط پوری کرن لئی ضروری نیں۔ سو ایس
تعلق دی تعریف ایویں وی ہو سکدی اے:

\begin{defn}[پابند متغیر دا ناں بدلنا، $\aconvone$] \ollabel{defn:aconvone}
  \emph{پابند متغیر دا ناں بدلنا} ($\aconvone$) استقرا نال
  ایویں متعین ہوندا اے۔
  \begin{enumerate}
  \item جے $N \aconvone N'$، تاں $\lambd[x][N] \aconvone
    \lambd[x][N']$۔ \ollabel{defn:aconvone1}
  \item جے $P \aconvone P'$، تاں $(PQ) \aconvone (P'Q)$۔ \ollabel{defn:aconvone2}
  \item جے $Q \aconvone Q'$، تاں $(PQ) \aconvone (PQ')$۔ \ollabel{defn:aconvone3}
  \item جے $x \neq y$، $y \notin \FV{N}$ تے $\Subst{N}{y}{x}$ معین اے، تاں
    $\lambd[x][N] \redone[\alpha] \lambd[y][\Subst{N}{y}{x}]$۔
    \ollabel{defn:aconvone4}
  \end{enumerate}
\end{defn}

دوویں تعریفاں ہم ارز نیں۔ ایہدا ثبوت اسیں مشق لئی چھڈدے آں۔
اگّوں اسیں استقرائی تعریف ورتاں گے۔

\begin{defn}[$\alpha$-روپ بدلنا، $\aconv$]
  \emph{$\alpha$-روپ بدلنا} ($\aconv$) ٹرماں اُتے سبھ توں چھوٹا
  انعکاسی تے تعدّی والا تعلق اے جیہدے اندر $\aconvone$ شامل اے۔
\end{defn}

اُتے وانگ، «سبھ توں چھوٹا» دا مطلب اے کہ تعلق وچ بس اوہی
جوڑے نیں جیہڑے تعدّی تے $\aconvone$ کر کے ضروری نیں۔ ایس توں
ہِٹھلی ہم ارز تعریف ملدی اے:

\begin{defn}[$\alpha$-روپ بدلنا، $\aconv$] \ollabel{defn:aconv}
  \emph{$\alpha$-روپ بدلنا} ($\aconv$) استقرا نال ایویں متعین ہوندا اے۔
  \begin{enumerate}
  \item جے $P \aconv Q$ تے $Q \aconv R$، تاں $P \aconv R$۔
    \ollabel{defn:aconv1}
  \item جے $P \aconvone Q$، تاں $P \aconv Q$۔ \ollabel{defn:aconv2}
  \item $P \aconv P$۔ \ollabel{defn:aconv3}
  \end{enumerate}
\end{defn}

\begin{ex}
  $\lambd[x][f x]$ دا $\alpha$-روپ بدل کے $\lambd[y][f
  y]$ ملدا اے، تے اُلٹا وی۔ سادہ لفظاں وچ، دوویں ایہو جیہے
  فنکشن نیں جیہڑے اک دلیل لیندے نیں تے اوہ دلیل اُتے $f$ لاؤن دا
  نتیجہ دیندے نیں۔ ایتھے حوالہ ماحول دے متغیر~$f$ دا اے۔

  $\lambd[x][f x]$ دا $\alpha$-روپ بدل کے $\lambd[x][g
    x]$ نہیں ملدا۔ سادہ لفظاں وچ، ایہ بالترتیب ماحول دے متغیراں
  $f$ تے~$g$ ول اشارہ کردے نیں، ایس لئی ایہ وکھرے فنکشن نیں۔
  جیہڑے ماحول وچ $f$ تے $g$ وکھرے ہون، اوتھے ایہناں دا ورتاؤ
  وی وکھرا اے۔
\end{ex}

\begin{prob}
  کی ہِٹھلے ٹرماں دے جوڑیاں وچ $\alpha$-روپ بدلیا جا سکدا اے؟
  \begin{enumerate}
  \item $\lambd[x][\lambd[y][x]]$ تے $\lambd[y][\lambd[x][y]]$
  \item $\lambd[x][\lambd[y][x]]$ تے $\lambd[c][\lambd[b][a]]$
  \item $\lambd[x][\lambd[y][x]]$ تے $\lambd[c][\lambd[b][a]]$
  \end{enumerate}
\end{prob}

\begin{lem}\ollabel{lem:fv-one}
  جے $P \aconvone Q$، تاں $\FV{P} = \FV{Q}$۔
\end{lem}

\begin{proof}
  $P \aconvone Q$ دے !!{derivation} اُتے استقرا نال۔
  \begin{enumerate}
  \item جے آخری قاعدہ \olref{defn:aconvone4} اے، تاں $P$ دی
    شکل $\lambd[x][N]$ تے $Q$ دی شکل
    $\lambd[y][\Subst{N}{y}{x}]$ اے، جتھے $x \neq y$، $y \notin \FV{N}$
    تے $\Subst{N}{y}{x}$ معین اے۔ اسیں $x \in \FV{N}$ دے
    ہون یا نہ ہون مطابق صورتاں وکھ کردے آں۔
    \begin{enumerate}
    \item جے $x \in FV(N)$، تاں۔
      \begin{align*}
        \FV{\lambd[y][\Subst{N}{y}{x}]} &=
        \FV{\Subst{N}{y}{x}} \setminus \{y\} \\
        & = ((\FV{N} \setminus \{x\}) \cup \{y\}) \setminus \{y\}
         && \text{ \olref[sub]{thm:infv} مطابق} \\
        & = \FV{N} \setminus \{x\} \\
        & = \FV{\lambd[x][N]}
      \end{align*}
    \item جے $x \notin FV(N)$، تاں۔
      \begin{align*}
        \FV{\lambd[y][\Subst{N}{y}{x}]}
        & = FV{\Subst{N}{y}{x}} \setminus \{y\} \\
        & = \FV{N} \setminus \{x\}
         && \text{\olref[sub]{thm:notinfv} مطابق} \\
        & = \FV{\lambd[x][N]}.
      \end{align*}
    \end{enumerate}
  \item باقی تِن صورتاں مشق لئی چھڈدے آں۔
  \end{enumerate}
\end{proof}

\begin{prob}
  \olref[lam][syn][alp]{lem:fv-one} دا ثبوت پورا کرو۔
\end{prob}

\begin{lem}\ollabel{lem:inv}
  جے $P \aconvone Q$، تاں $Q \aconvone P$۔
\end{lem}

\begin{proof}
  $P \aconvone Q$ دے !!{derivation} اُتے استقرا۔
  \begin{enumerate}
  \item جے آخری قاعدہ \olref{defn:aconvone4} اے، تاں $P$ دی
    شکل $\lambd[x][N]$ تے $Q$ دی شکل
    $\lambd[y][\Subst{N}{y}{x}]$ اے، جتھے $x \neq y$، $y \notin \FV{N}$
    تے $\Subst{N}{y}{x}$ معین اے۔ پہلاں، $y \notin
    \FV{\Subst{N}{y}{x}}$، \olref[sub]{thm:clr} مطابق۔
    \olref[sub]{thm:inv} توں ملدا اے کہ $\Subst{\Subst{N}{y}{x}}{x}{y}$
    معین اے تے~$N$ دے برابر وی اے۔ فیر \olref{defn:aconvone4}
    مطابق $\lambd[y][\Subst{N}{y}{x}]
    \aconvone \lambd[x][\Subst{\Subst{N}{y}{x}}{x}{y}] =
    \lambd[x][N]$۔
  \end{enumerate}
\end{proof}

\begin{prob}
  \olref[lam][syn][alp]{lem:inv} دا ثبوت پورا کرو۔
\end{prob}

\begin{thm}
  $\alpha$-روپ بدلنا ٹرماں اُتے اک ہم ارزیتی تعلق اے، یعنی
  ایہ انعکاسی، متقابل تے تعدّی والا اے۔
\end{thm}

\begin{proof}
  \begin{enumerate}
  \item ہر ٹرم $M$ لئی، $M$ نوں $M$ وچ \emph{صفر} پابند ناں بدل کے
    بدلیا جا سکدا اے۔
  \item جے $P$ دا $\alpha$-روپ پابند ناں بدلن دے اک سلسلے نال
    $Q$ بنے، تاں $Q$ توں اسیں ایہناں بدلاں نوں
    \olref{lem:inv} مطابق اُلٹی ترتیب وچ واپس کر کے~$P$ حاصل
    کر سکدے آں۔
  \item جے $P$ دا $\alpha$-روپ پابند ناں بدلن دے اک سلسلے نال
    $Q$، تے $Q$ دا دوجے سلسلے نال $R$ بنے، تاں پہلا سلسلہ
    تے فیر دوجا سلسلہ ورت کے $P$ نوں $R$ وچ بدل سکدے آں۔
  \end{enumerate}
\end{proof}

اگّوں اسیں آکھاں گے کہ $M$ تے $N$ \emph{$\alpha$-ہم ارز}
نیں، $M
\aeq N$، جے تے صرف جے $M$ دا $\alpha$-روپ بدل کے $N$ ملے۔
جیویں ہُن دکھایا اے، ایہ اُدوں تے صرف اُدوں ہوندا اے جدوں $N$
دا $\alpha$-روپ بدل کے~$M$ ملے۔

\begin{thm}\ollabel{thm:fv}
  جے $M \aeq N$، تاں $\FV{M} = \FV{N}$۔
\end{thm}

\begin{proof}
  \olref{lem:fv-one} توں سیدھا ملدا اے۔
\end{proof}

\begin{lem}\ollabel{lem:sub:R}
  جے $R \aeq R'$ تے $\Subst{M}{R}{y}$ معین اے، تاں
  $\Subst{M}{R'}{y}$ معین اے تے $\alpha$-ہم ارز اے
  $\Subst{M}{R}{y}$ نال۔
\end{lem}

\begin{proof}
  مشق۔
\end{proof}

\begin{prob}
  \olref[lam][syn][alp]{lem:sub:R} ثابت کرو۔
\end{prob}

یاد کرو کہ \olref[sub]{sec} وچ کجھ صورتاں لئی بدل غیر معین اے۔
پر ٹرماں دا $\alpha$-روپ بدل کے اسیں پابند متغیراں دے ناں
ایویں بدل سکدے آں کہ بدل ہمیشہ معین ہووے۔ نتیجہ $\alpha$-ہم ارزی
قائم رکھدا اے، جیویں ایہ قضیہ دسدا اے:

\begin{thm}\ollabel{thm:sub}
  ہر $M$، $R$ تے~$y$ لئی، کوئی $M'$ موجود اے جیہدے لئی $M \aeq M'$
  تے $\Subst{M'}{R}{y}$ معین اے۔ نالے جے کوئی ہور جوڑا
  $M'' \aeq M$ تے $R''$ اے جیہدے لئی $\Subst{M''}{R''}{y}$
  معین اے تے $R'' \aeq R$، تاں $\Subst{M'}{R}{y} \aeq \Subst{M''}{R''}{y}$۔
\end{thm}

\begin{proof}
  $M$ دی بناوٹ اُتے استقرا نال۔
  \begin{enumerate}
  \item $M$ اک متغیر~$z$ اے۔ مشق۔
  \item فرض کرو $M$ دی شکل $\lambd[x][N]$ اے۔ ایہو جیہا متغیر
    $z$ چُنو جیہڑا $x$ تے $y$ توں وکھرا ہووے تے جیہدے لئی
    $z \notin \FV{N}$ تے $z
    \notin \FV{R}$۔ استقرائی مفروضے مطابق کوئی $N'$ اے جیہدے
    لئی $N' \aeq N$ تے $\Subst{N'}{z}{x}$ معین اے۔ فیر
    $\lambd[x][N] \aeq \lambd[x][N']$ وی اے، جیویں
    \olref{defn:aconvone}\olref{defn:aconvone1} توں ملدا اے۔
    ہُن $\lambd[x][N']
    \aeq \lambd[z][\Subst{N'}{z}{x}]$ اے،
    \olref{defn:aconvone}\olref{defn:aconvone4} مطابق۔ ایہ
    کر سکدے آں کیوں جے $z \ne x$، $z \notin FV(N')$ تے
    $\Subst{N'}{z}{x}$ معین اے۔ آخر وچ
    $\Subst{\lambd[z][\Subst{N'}{z}{x}]}{R}{y}$ معین اے،
    کیوں جے $z \neq y$ تے $z \notin FV(R)$۔

    نالے جے کوئی ہور $N''$ تے $R''$ اوہو شرطاں پوریاں کرن، تاں۔
    \begin{multline*}
      \Subst{(\lambd[z][\Subst{N''}{z}{x}])}{R''}{y} =\\
    \begin{aligned}
      &= \lambd[z][\Subst{\Subst{N''}{z}{x}}{R''}{y}] \\
      &= \lambd[z][\Subst{\Subst{N''}{z}{x}}{R}{y}] && \text{
        \olref{lem:sub:R} مطابق}\\
      &=\lambd[z][\Subst{\Subst{N'}{z}{x}}{R}{y}]
      && \text{استقرائی
        مفروضے مطابق}\\
      &=\Subst{(\lambd[z][\Subst{N'}{z}{x}])}{R}{y}
    \end{aligned}
    \end{multline*}
    \item $M$ دی شکل $(PQ)$ اے۔ مشق۔
  \end{enumerate}
\end{proof}

\begin{prob}
  \olref[lam][syn][alp]{thm:sub} دا ثبوت پورا کرو۔
\end{prob}

\begin{cor}\ollabel{cor:sub}
  ہر $M$، $R$ تے~$y$ لئی $M'$ تے $R'$ دا اک جوڑا موجود اے
  جیہدے لئی $M' \aeq M$، $R \aeq R'$ تے $\Subst{M'}{R'}{y}$
  معین اے۔ نالے جے کوئی ہور جوڑا $M'' \aeq M$ تے $R''$ اے
  جیہدے نال $\Subst{M'}{R'}{y}$ معین اے، تاں $\Subst{M'}{R'}{y} \aeq
  \Subst{M''}{R''}{y}$۔
\end{cor}

\begin{proof}
  \olref{thm:sub} توں سیدھا ملدا اے۔
\end{proof}

\end{document}
